\reportheading{附录 F.\quad 两个代表工况的 FDS 输入}
丙烷算例取 OP-6、106.2~kW 的粗网格验证工况。丙烯算例取推荐风量 1.12~m$^3$/s 下控制净高的 OP-6 工况，烟炱产率为 0.05。两组均使用 FDS 6.10.1，环境温度取 $27.5\,^{\circ}$C，默认壁面为绝热。以下坐标均从模型原点起算，单位为 m。

\tabletitle{表 F1. 与输入文件和求解记录核对的主要设定。}
\begin{reporttable}{L{.22}L{.38}L{.40}}
\rowcolor{tablegrey}项目 & 丙烷验证 & 丙烯排烟设计\\
网格与辐射 & $132\times28\times19$，辐射份额 0.29 & $132\times28\times38$，辐射份额 0.37\\
燃烧 & \texttt{PROPANE}，名义功率 106.2~kW & \texttt{PROPYLENE}，名义功率 106.2~kW。烟炱和 CO 产率分别为 0.05 和 0.01~kg/kg 燃料。\\
端墙开口 & 房间两端封闭 & $x=0$ 和 6.60 处设 \texttt{OPEN}，范围为 $y=0$--0.90、$z=0$--0.48\\
挡烟垂壁 & 无 & 位于 $x=2.80$ 和 3.80，范围为 $y=0$--0.90、$z=0.52$--0.76\\
顶棚排烟口 & 无 & 左口 $x=2.90$--3.30，右口 $x=3.30$--3.70。两口均在 $z=0.76$，宽度方向为 $y=0.05$--0.85。\\
测点与分析时段 & 51 支直径 1~mm 的 \texttt{THERMOCOUPLE}，位于 $y=0.45$、$z=0.74$。温升取 90--120~s 平均。 & 在 $y=0.45$ 的八个测站设 $z=0$--0.76 的 \texttt{LAYER HEIGHT} 测柱（第 2.2 节）。同一测站在 $z=0.38$ 记录 CO 和能见度。\\
\end{reporttable}

\subhead{两组共用的边界与燃烧器}
计算域为 $x=0$--6.60、$y=0$--1.40、$z=0$--0.76。侧墙占 $y=0.90$--0.95，R1 门洞沿 $x=3.91$--4.34、$z=0$--0.63 从墙中切出。门外延伸域在 $y=1.40$ 的外侧面，以及 $y=0.95$--1.40 的两端和顶部设 \texttt{OPEN} 边界。两组气相均用 Deardorff 模型，计算至 120~s，设备与热释放率的名义输出间隔均为 0.5~s。

燃烧面贴地，输入坐标为 $x=3.175$--3.425、$y=0.325$--0.575。对齐网格后，两组均占 $5\times5$ 个边长 0.05~m 的水平单元，面积为 0.0625~m$^2$。火源面设 \texttt{HRRPUA} = 1699.2~kW/m$^2$，启动时采用 FDS 默认因子 $\tanh(t/1\,\mathrm{s})$~\cite{fdsuser}。

\subhead{排烟设定与记录流量}
每个排烟口通过一条 HVAC 支管接至环境节点。每条支管截面积为 0.32~m$^2$，长度为 0.40~m，\texttt{VOLUME\_FLOW} 设为 0.56~m$^3$/s。该工况两条支管的 \texttt{DUCT VOLUME FLOW} 记录在 0--120~s 内的合计均值为 1.11345~m$^3$/s。表 C3 给出设定风量 1.12~m$^3$/s 下的净高结果。

\clearpage
\reportheading{附录 G.\quad PINN 训练与最低净高校正}
第 5 节的校正 PINN 先拟合八个 FDS 测站的烟层高度时序，再逐站校正预测的最低净高。模型输入包括时间和七个工况参数。表 G1 的几何与流量参数分别以 m 和 m$^3$/s 为单位。

\tabletitle{表 G1. 第 5 节代理模型所用工况输入。}
\begin{reporttable}{L{.45}L{.55}}
\rowcolor{tablegrey}输入 & 送入模型的数值\\
左右端墙补风口高度 $a_L,a_R$ & $(a_L-0.44)/0.08$、$(a_R-0.44)/0.08$\\
门态 & OP-1 取 0，OP-6 取 1\\
烟炱产率 $Y_s$ & $(Y_s-0.05)/0.05$\\
挡烟垂壁深度 $d$ & $(d-0.24)/0.08$\\
总排烟量 $Q$ & $(Q-1.28)/0.32$\\
竖向网格 & 19 层取 0，38 层取 1\\
\end{reporttable}

\subhead{时序拟合与上层体积平衡}
时间基函数包含常数项和 96 个固定的 tanh 项。各 tanh 项的自变量为 $(t-c)/s+\boldsymbol b^\mathsf{T}\boldsymbol u$，$\boldsymbol u$ 为工况输入向量。中心 $c$ 在 0--120~s 内均匀抽取，尺度 $s$ 在 1--25~s 内按对数均匀分布抽取，输入权重 $\boldsymbol b$ 服从均值为 0、标准差为 0.35 的正态分布。时间参数和输入权重分别由两个 NumPy 随机数生成器抽取，两者的种子均为 3322。将基函数乘以 $1-\exp[-t/(2\,\mathrm{s})]$，使初始烟层界面高度为 $H=0.76$~m。输出系数按岭惩罚系数 0.1 拟合 FDS 上层厚度 $H-h_j$。

设 $\bar d(t)$ 为预测上层厚度的空间平均值，权重取各测站对应的纵向区段长度，区段以相邻测站中点和端墙为界。拟合采用以下上层体积平衡式。
\[
 r(t)=\frac{\mathrm d\bar d}{\mathrm dt}
       +\frac{Q}{(6.60\,\mathrm m)(0.90\,\mathrm m)H}\,\bar d(t),
 \qquad r(t)\approx \overline{r}_{t\geq10\,\mathrm{s}} .
\]
对每个工况减去 $t\geq10$~s 的平均残差，以消去未知的常数净源项。拟合在单位欧氏范数的空间平均方向上施加残差平方惩罚，权重为 25。各测站之间的差异由 FDS 数据拟合。

\subhead{逐测站最低值校正}
先将各测站的预测曲线限制在 $[0,H]$，再在 0--120~s 内每隔 0.5~s 取值，求得预测最低值。对每次训练用的 FDS 计算，校正目标取 FDS 实际最低值减去该站的拟合最低值。以高斯核 $k(\boldsymbol u,\boldsymbol v)=\exp(-\|\boldsymbol u-\boldsymbol v\|^2/2)$ 拟合中心化后的校正量，岭惩罚系数取 0.25，预测后加回训练均值。各测站的预测最低值加上校正量后再次限制在 $[0,H]$，其中最低值作为该工况的 PINN 预测。每次留出均将待测设计及其网格变体同时排除在时序拟合和校正之外。直接回归用同一核函数和惩罚系数拟合各工况的整体最低值，混合模型对两种预测取平均。
